% Part: first-order-logic
% Chapter: natural-deduction
% Section: provability-quantifiers

% verification of properties of provability needed for maximally
% consistent sets in the completeness chapter.

\documentclass[../../../include/open-logic-section]{subfiles}

\begin{document}

\olfileid{fol}{ntd}{qpr}

\olsection{可推导性与量词}

\begin{explain}
完备性定理还要求自然演绎规则能给出本节所建立的关于~$\Proves$ 的结论。
\end{explain}

\begin{thm}
\ollabel{thm:strong-generalization} 若常项符号~$c$ 不在~$\Gamma$ 或~$!A(x)$ 中出现，且~$\Gamma \Proves !A(c)$，则 $\Gamma
\Proves \lforall[x][!A(x)]$。
\end{thm}

\begin{proof}
设~$\delta$ 是从~$\Gamma$ 推出~$!A(c)$ 的一个推导。通过添加一个~$\Intro{\lforall}$ 推理步骤，我们得到从~$\Gamma$ 推出 $\lforall[x][!A(x)]$ 的一个推导。由于~$c$ 不在~$\Gamma$ 或~$!A(x)$ 中出现，故满足特征变元条件。
\end{proof}

\begin{prop}
\ollabel{prop:provability-quantifiers}
\begin{tagenumerate}{prvEx,prvAll}
\tagitem{prvEx}{$!A(t) \Proves \lexists[x][!A(x)]$.}{}

\tagitem{prvAll}{$\lforall[x][!A(x)] \Proves !A(t)$.}{}
\end{tagenumerate}
\end{prop}

\begin{proof}
\begin{tagenumerate}{prvEx,prvAll}
\tagitem{prvEx}{以下是从~$!A(t)$ 推出~$\lexists[x][!A(x)]$ 的推导：
\begin{prooftree}
\AxiomC{$!A(t)$}
\RightLabel{$\Intro{\lexists}$}
\UnaryInfC{$\lexists[x][!A(x)]$}
\end{prooftree}}{}

\tagitem{prvAll}{以下是从~$\lforall[x][!A(x)]$ 推出~$!A(t)$ 的推导：
\begin{prooftree}
\AxiomC{$\lforall[x][!A(x)]$}
\RightLabel{$\Elim{\lforall}$}
\UnaryInfC{$!A(t)$}
\end{prooftree}}{}
\end{tagenumerate}
\end{proof}

\end{document}